\documentclass[12pt]{article}
\usepackage[utf8]{inputenc}
\usepackage[T1]{fontenc}
\usepackage[margin=1in]{geometry}
\usepackage{setspace}
\usepackage{xcolor}
\usepackage[colorlinks=true, urlcolor=blue!60!black, linkcolor=blue!60!black]{hyperref}
\setlength{\parindent}{0em}
\setlength{\parskip}{0.6em}
\renewcommand{\baselinestretch}{1.0}
\begin{document}
\begin{center}
{\LARGE \textit{Bivens v.\ Six Unknown Named Agents of Federal Bureau of Narcotics}, 403 U.S.\ 388 (1971)\par}
\vspace{0.5em}
{\large Case Brief --- Two-Page Edition\par}
\end{center}

\section{Memory Jogger}
Victims of unconstitutional searches by federal agents may sue for damages directly under the Fourth Amendment (\href{https://supreme.justia.com/cases/federal/us/403/388/}{Justia}).

\section{Detailed Case Facts}
Federal narcotics agents entered Webster Bivens's Brooklyn apartment without a warrant or probable cause, manacled him in front of his family, searched the apartment ``from stem to stern,'' and strip-searched him at the courthouse.

\section{Procedural History}
No statute authorized damages against federal officers, since \S~1983 reaches only state officials. The district court dismissed, and the Second Circuit affirmed.

\section{Judicial Votes}
Reversed, 6--3. Brennan wrote for Douglas, Stewart, White, and Marshall; Harlan concurred in the judgment; Burger, Black, and Blackmun dissented.

\section{Holding}
A Fourth Amendment violation by federal agents supports a damages action under the Constitution itself, absent special factors counseling hesitation or an equally effective remedy chosen by Congress.

\section{Analysis of Opinions}
\textbf{Majority:} Federal power ``does not disappear like a magic gift when it is wrongfully used''; state tort law may be hostile to Fourth Amendment interests; and damages are the traditional remedy for invasions of liberty. \textbf{Harlan:} courts that may enjoin may award damages, and the exclusionary rule does nothing for an innocent victim. \textbf{Dissents:} creating remedies is for Congress, and the Court invited an ``avalanche'' of suits. \textbf{Later:} \textit{Abbasi}, \textit{Egbert}, and \textit{Goldey} have confined \textit{Bivens} to its own contexts.

\section{Examples: Future Applications}
\textbf{Same-Side:} a warrantless, abusive DEA home raid. \textbf{Opposite-Side:} a TSA screener's excessive force, a new context after \textit{Egbert}. \textbf{Fence-Sitter:} a \textit{Bivens}-like raid by FBI agents in a fraud case.

\section{Critique}
\textit{Bivens} pits \textit{Marbury}'s right-remedy maxim against separation of powers, and its narrowing leaves federal agents less accountable than state police under \S~1983.

\section{Key Quotations}
\begin{quote}
``The present case involves no special factors counseling hesitation in the absence of affirmative action by Congress. ... For we have here no explicit congressional declaration that persons injured by a federal officer's violation of the Fourth Amendment may not recover money damages from the agents, but must instead be remitted to another remedy, equally effective in the view of Congress.''
\end{quote}
\noindent --- Justice Brennan, for the Court, 403 U.S.\ at 396--97.

\end{document}
